\documentclass[12pt,twoside]{article}

\usepackage[english]{babel}
\usepackage[utf8]{inputenc}
\usepackage{sansmathfonts}
\usepackage[T1]{fontenc}
\renewcommand*\familydefault{\sfdefault} %% Only if the base font of the document is to be sans serif
%\usepackage{tabularx}
\usepackage{dashrule}
\usepackage{tikz}
\usepackage{tabularray}
\usepackage[export]{adjustbox}
%\captionsetup[longtable]{skip=30pt}
\usepackage{arydshln}
\usepackage{ctable}
\usepackage{colortbl}
\usepackage{multicol,lipsum}
\usepackage{amsmath}
\usepackage{amssymb,amsthm,amsmath,dsfont,mathrsfs}
\usepackage[colorinlistoftodos]{todonotes}
\usepackage{geometry}
\geometry{
a4paper,
total={170mm,257mm},
left=15mm,
right=15mm,
top=15mm,
bottom=10mm,
}
\usepackage{natbib}
\usepackage{soul}
\usepackage{setspace}
\usepackage{hyperref}
\usepackage{xcolor}
\usepackage{cancel}
\color{darkgray}
\usepackage{graphicx}
\usepackage{fancyhdr}
\pagestyle{fancy}
\renewcommand{\headrulewidth}{0pt}


\newcommand{\T}[1]{\textbf{\large #1} \newline}

\newcommand{\TM}[1]{ \textbf{ \newline #1}  }
\setlength\columnsep{30pt}

\definecolor{bluebell}{rgb}{0.39, 0.44, 1}
\definecolor{laranja}{rgb}{1, 0.58, 0.33}
\definecolor{verde}{rgb}{0, 0.46, 0.37}
\definecolor{babypink}{rgb}{0.85, 0.65, 0.65} 
\definecolor{cinza}{rgb}{0.3, 0.3, 0.3} 
\definecolor{iris}{rgb}{0.39, 0.44, 1} 
\definecolor{babyblue}{rgb}{0.13, 0.59, 0.95} 

\color{cinza} 
 
\newcommand{\C}[1]{\textcolor{cinza}{#1}}
\newcommand{\BI}[1]{\textbf{\textcolor{iris}{#1}}}
\newcommand{\BL}[1]{\textbf{\textcolor{laranja}{#1}}}
\newcommand{\BV}[1]{\textbf{\textcolor{verde}{#1}}}
\newcommand{\BC}[1]{\textbf{\textcolor{cinza}{#1}}}
\newcommand{\BB}[1]{\textbf{\textcolor{babyblue}{#1}}}

\begin{document}
\pagenumbering{gobble}
\setstretch{1}

\setstretch{1} 
\normalsize 
\setlength{\tabcolsep}{0em} 
\setlength\parindent{0pt} 
\begin{tabular}{p{0.19\textwidth} p{0.61\textwidth} p{0.2\textwidth}}
\BB{\Large Exercícios } & \footnotesize \BI{Módulo Avançado} \BL{-} \BV{Probabilidade} \BL{-} \BC{Nível I} & Nota:    \\  
\end{tabular}
\vspace{0.3cm}

\small
\textcolor{laranja}{Instruções: } \textcolor{cinza}{Resolva os seguintes exercícios.}
\vspace{0.1cm}

\setstretch{1.3}
\footnotesize
\begin{tblr}{
colspec={p{0.46\textwidth}|[2pt,laranja]p{0.46\textwidth}},
columns = {colsep=0pt, rightsep=10pt},
column{2} = {leftsep=8pt},
rows={valign=h, ht=8.200000000000001cm},
}
\BI{1)~}Uma bolsa contém 10 conchas grandes e 17 pequenas. Qual a chance de tirar uma pequena?\vspace{5px}  \newline &\BI{2)~}Em um sorteio com 56 prêmios dourados e prateados, a chance de ganhar um prêmio dourado é \(\dfrac{9}{14}\). Quantos prêmios dourados há?\vspace{5px}  \newline \\ 
\BI{3)~}Num pote há 14 clips de metal, 17 clips de plástico e 14 clips de papel. Qual a chance de pegar um clip de metal?\vspace{5px}  \newline &\BI{4)~}Em uma caixa há 8 moedas de R$1 e outras de R$0,50. Se a chance de tirar uma de R$1 é \(\dfrac{4}{21}\), quantas moedas há?\vspace{5px}  \newline \\ 
\BI{5)~}Há 72 alunos numa sala, sendo homens e mulheres. Se a probabilidade de escolher um homem é \(\dfrac{1}{2}\), quantos homens há?\vspace{5px}  \newline &\BI{6)~}Há algumas frutas maduras e 16 verdes. Se a chance de escolher uma verde é \(\dfrac{4}{5}\), qual o total de frutas no cesto?\vspace{5px}  \newline \\ 
\end{tblr}

\newpage 
\setstretch{1} 
\normalsize 
\setlength{\tabcolsep}{0em} 
\setlength\parindent{0pt} 
\BB{\Large Resolução} \quad \footnotesize \BI{Módulo Avançado} \BL{-} \BV{Probabilidade} \BL{-} \BC{Nível I}
\vspace{0.3cm}

\small
\textcolor{laranja}{Instruções: } \textcolor{cinza}{Resolva os seguintes exercícios.}
\vspace{0.1cm}



\setstretch{1.5}
\footnotesize
\begin{tblr}{
colspec={p{0.46\textwidth}|[2pt,laranja]p{0.46\textwidth}},
columns = {colsep=0pt, rightsep=10pt},
column{2} = {leftsep=8pt},
rows={valign=h, ht=8.200000000000001cm},
}
\BI{1)~}{Uma bolsa contém 10 conchas grandes e 17 pequenas. Qual a chance de tirar uma pequena?}\vspace{5px}  \newline 
{$\definecolor{azulEscuro}{rgb}{0.25, 0.35, 0.89}\definecolor{rosaClaro}{rgb}{1, 0.52, 0.60}\textcolor{azulEscuro}{P(A)=\dfrac{n(A)}{n(\Omega)}} \vspace{5px} \\A: \text{Conchas pequenas} \\n(A)=17,\quad n(\Omega)=10+17=27 \vspace{5px} \\P(A)=\dfrac{17}{27}$} &\BI{2)~}{Em um sorteio com 56 prêmios dourados e prateados, a chance de ganhar um prêmio dourado é \(\dfrac{9}{14}\). Quantos prêmios dourados há?}\vspace{5px}  \newline 
{$\definecolor{azulEscuro}{rgb}{0.25, 0.35, 0.89}\definecolor{rosaClaro}{rgb}{1, 0.52, 0.60}\textcolor{azulEscuro}{P(A)=\dfrac{n(A)}{n(\Omega)}} \vspace{5px} \\A: \text{Prêmios dourados} \vspace{5px} \\P(A)=\dfrac{9}{14},\quad n(\Omega)=56 \vspace{5px} \\\dfrac{9}{14}=\dfrac{n(A)}{56} \vspace{5px} \\n(A)=\dfrac{9}{14}\cdot 56=\dfrac{504}{14}=36$} \\ 
\BI{3)~}{Num pote há 14 clips de metal, 17 clips de plástico e 14 clips de papel. Qual a chance de pegar um clip de metal?}\vspace{5px}  \newline 
{$\definecolor{azulEscuro}{rgb}{0.25, 0.35, 0.89}\definecolor{rosaClaro}{rgb}{1, 0.52, 0.60}\textcolor{azulEscuro}{P(A)=\dfrac{n(A)}{n(\Omega)}} \vspace{5px} \\A: \text{Clips de metal} \\n(A)=14,\quad n(\Omega)=14+17+14=45 \vspace{5px} \\P(A)=\dfrac{14}{45}$} &\BI{4)~}{Em uma caixa há 8 moedas de R$1 e outras de R$0,50. Se a chance de tirar uma de R$1 é \(\dfrac{4}{21}\), quantas moedas há?}\vspace{5px}  \newline 
{$\definecolor{azulEscuro}{rgb}{0.25, 0.35, 0.89}\definecolor{rosaClaro}{rgb}{1, 0.52, 0.60}\textcolor{azulEscuro}{P(A)=\dfrac{n(A)}{n(\Omega)}} \vspace{5px} \\A: \text{Moedas de R$1} \vspace{5px} \\P(A)=\dfrac{4}{21},\quad n(A)=8 \vspace{5px} \\\dfrac{4}{21}=\dfrac{8}{n(\Omega)} \vspace{5px} \\n(\Omega)=\dfrac{8}{\dfrac{4}{21}} \vspace{5px} \\n(\Omega)=8\cdot \dfrac{21}{4}=\dfrac{168}{4}=42$} \\ 
\BI{5)~}{Há 72 alunos numa sala, sendo homens e mulheres. Se a probabilidade de escolher um homem é \(\dfrac{1}{2}\), quantos homens há?}\vspace{5px}  \newline 
{$\definecolor{azulEscuro}{rgb}{0.25, 0.35, 0.89}\definecolor{rosaClaro}{rgb}{1, 0.52, 0.60}\textcolor{azulEscuro}{P(A)=\dfrac{n(A)}{n(\Omega)}} \vspace{5px} \\A: \text{Homens} \vspace{5px} \\P(A)=\dfrac{1}{2},\quad n(\Omega)=72 \vspace{5px} \\\dfrac{1}{2}=\dfrac{n(A)}{72} \vspace{5px} \\n(A)=\dfrac{1}{2}\cdot 72=\dfrac{72}{2}=36$} &\BI{6)~}{Há algumas frutas maduras e 16 verdes. Se a chance de escolher uma verde é \(\dfrac{4}{5}\), qual o total de frutas no cesto?}\vspace{5px}  \newline 
{$\definecolor{azulEscuro}{rgb}{0.25, 0.35, 0.89}\definecolor{rosaClaro}{rgb}{1, 0.52, 0.60}\textcolor{azulEscuro}{P(A)=\dfrac{n(A)}{n(\Omega)}} \vspace{5px} \\A: \text{Frutas verdes} \vspace{5px} \\P(A)=\dfrac{4}{5},\quad n(A)=16 \vspace{5px} \\\dfrac{4}{5}=\dfrac{16}{n(\Omega)} \vspace{5px} \\n(\Omega)=\dfrac{16}{\dfrac{4}{5}} \vspace{5px} \\n(\Omega)=16\cdot \dfrac{5}{4}=\dfrac{80}{4}=20$} \\ 
\end{tblr}

\end{document}
